\documentclass[11pt,a4paper]{article}
\usepackage[T1]{fontenc}
\usepackage{lmodern}
\usepackage[margin=27mm,headheight=14pt]{geometry}
\usepackage{amsmath,amssymb,amsthm,mathtools}
\usepackage{microtype,booktabs,array,longtable}
\usepackage{xcolor}
\usepackage{enumitem}
\usepackage{fancyhdr}
\usepackage[hyphens]{url}
\usepackage[colorlinks=true,linkcolor=blue!45!black,citecolor=blue!45!black,urlcolor=blue!45!black]{hyperref}
\usepackage[nameinlink,noabbrev]{cleveref}
\setlength{\emergencystretch}{2em}
\setlength{\parskip}{4pt}
\setlength{\parindent}{1.1em}
\setlist{itemsep=3pt,topsep=5pt}
\pagestyle{fancy}
\fancyhf{}
\fancyhead[L]{\small Exponential-feedback transseries}
\fancyhead[R]{\small Sharp large-order asymptotics}
\fancyfoot[C]{\thepage}
\newtheorem{theorem}{Theorem}[section]
\newtheorem{proposition}[theorem]{Proposition}
\newtheorem{lemma}[theorem]{Lemma}
\newtheorem{corollary}[theorem]{Corollary}
\theoremstyle{definition}
\newtheorem{definition}[theorem]{Definition}
\newtheorem{remark}[theorem]{Remark}
\newtheorem{example}[theorem]{Example}
\newcommand{\N}{\mathbb N}
\newcommand{\R}{\mathbb R}
\newcommand{\Z}{\mathbb Z}
\newcommand{\E}{\mathbb E}
\newcommand{\Prob}{\mathbb P}
\newcommand{\TV}{d_{\mathrm{TV}}}
\newcommand{\Pois}{\operatorname{Pois}}
\newcommand{\U}{\widehat U}
\newcommand{\calM}{\mathcal M}
\newcommand{\dd}{\,\mathrm d}
\newcommand{\epsn}{\varepsilon_n}
\DeclareMathOperator{\supp}{supp}
\DeclareMathOperator{\Var}{Var}
\numberwithin{equation}{section}
\hypersetup{pdftitle={Microscopic Condensation in Exponential-Feedback Transseries},pdfauthor={Research draft prepared for Vladimir Reshetnikov},pdfsubject={Sharp coefficients, Gaussian-Poisson limits, finite-perturbation universality, and Borel growth}}
\title{\textbf{Microscopic Condensation in\\Exponential-Feedback Transseries}\\[7pt]
\large Sharp coefficients, Gaussian--Poisson limits,\\finite-perturbation universality, and refined Borel growth}
\author{Research draft prepared for Vladimir Reshetnikov}
\date{29 September 2026}
% ed. (2026-09-29): editorial-note environment for ProveIt cross-references;
% it uses no counter, so no author numbering changes.
\newenvironment{ednote}{\par\smallskip\begin{quote}\small\raggedright\textbf{Editorial note (ProveIt, 2026-09-29).}\ }{\end{quote}\smallskip}
\begin{document}
\maketitle
\begin{abstract}
Consider the formal fixed point
\(\U(q)=\sum_{j\ge1}q^j\exp(\lambda_j\U(q))\), with nonnegative slopes.
The ProveIt research article on this model leaves open a full multiplicative
coefficient equivalent and correctly centered fluctuation laws beyond its
logarithmic asymptotic. We resolve that question for an asymptotically quadratic
class, including every finite perturbation of \(\lambda_j=j^2\).
If
\(\lambda_j=j^2+\alpha j+\eta j/\log j+o(j/\log j)\),
\(b=1+\lambda_1\), and \(r_n(r_n+1)e^{2r_n}=n\), then
\[
 [q^n]\U(q)\sim
 \frac{\exp\!\left((2r_n+1)\frac{nr_n}{r_n+1}
                 +br_n^2+\alpha r_n+\eta/2\right)}
      {\sqrt{2r_n^2+4r_n+1}}.
\]
The proof identifies the full microscopic configuration: one large action,
a Poisson number of action-two factors, and otherwise action-one factors.
The large action has Gaussian fluctuations, asymptotically independent of the
Poisson cloud. This yields a three-level universality law, a multiplicative
positive-axis Borel-growth equivalent, and an equivalent for the least term.
In particular, changing only the first slope changes coefficients by a
lognormal-size factor, although it leaves the previously known logarithmic
asymptotic and Gevrey classification unchanged. All main results are proved
here by positive Lagrange weights and two discrete saddle analyses. Exact
finite computations accompany, but do not replace, the proofs.
\end{abstract}
% ed. (2026-09-29): the abstract's "We resolve that question" is stale at filing: the finite-core package resolved it first for pure power slopes.
\begin{ednote}
This article was written at a snapshot that did not yet contain the
package \nolinkurl{Finite_Core_Universality_Exponential_Feedback/}~\cite{ed:finitecore} in
\nolinkurl{Analysis/Transseries/docs/series-and-transseries/}, filed
before it (batch 47), and it does not cite that package. For slopes that
are eventually exactly $j^2$, \cref{thm:coefficient} is the $p=2$, $a=1$
case of that package's first-correction formula (its
Proposition~\texttt{prop:first-correction}, with $\lambda_1+w_2=b$), the variance
$2n/(\log n)^2$ is its equation \texttt{eq:variance}, and the least-term index of
\cref{thm:least} is its Corollary~\texttt{cor:least-term}; that package also
proves a local Gaussian law and a central limit theorem for the largest
action (its Theorem~\texttt{thm:clt}). The two are independent proofs of the same
coefficient theorem, and their constants agree. New here are the $\alpha j$
and $\eta j/\log j$ terms, the rate \eqref{eq:rate}, the joint
total-variation Gaussian--Poisson law of the largest part and the number of
twos, which answers that package's question on joint laws for the finite
background cloud at $p=2$, $a=1$, the two-dimensional local limit, the
Borel equivalent and the value of the least term. The total-variation
Poisson law at diverging intensity is also what the later package
\nolinkurl{Poisson_Layers_Finite_Core_Boundary/} lists as open; this
article proves it for this class. Neither package has been independently
reviewed.
\end{ednote}
\noindent\textbf{Scope and status.}
This is a mathematical research draft, not a Lean verification or a claim of
settling a named longstanding conjecture. It strengthens the specific
repository article identified in \cref{sec:context}. General publication
priority has not been established. The least-term result is not asserted to
be the remainder of an analytically summed solution.
\clearpage
\setcounter{tocdepth}{1}
\tableofcontents
\clearpage

\section{Repository context and the question being resolved}
\label{sec:context}
The relevant model is developed in ProveIt's
\emph{A Sharp Regularity Classification for Countable Exponential-Feedback
Transseries} \cite{repo-classification}. The inspected repository snapshot is
\begin{center}
\texttt{04473354a0f3edff2d6c365caf26170ca6b88735}.
\end{center}
The article lies under
\begin{quote}\small\ttfamily
Analysis/Transseries/docs/series-and-transseries/\\
Exponential\_Feedback\_Regularity\_Classification/article.tex.
\end{quote}
Its subsection \emph{Beyond logarithmic coefficient asymptotics} explicitly
asks for sublinear terms, a full multiplicative equivalent, and a local limit
law for the number of factors and the largest action. It points out that
small nonunit actions and the displacement of the large-action saddle become
relevant at that precision. Those are exactly the effects analyzed here.

For the quadratic model the previous logarithmic result is
\begin{equation}
 \log u_n=n\log n-2n\log\log n+(\log4-1)n+o(n).
 \label{eq:old-log}
\end{equation}
The previous concentration results locate the largest action at
\(2n/\log n\), up to a relative error tending to zero, and the previous Borel
result is \(\log\log B(t)\sim2\sqrt t\).
None of these assertions determines a multiplicative equivalent. In
particular, the factor \(e^{2r_n^2}\) obtained below is invisible in
\eqref{eq:old-log}, yet grows faster than every fixed power of \(n\).

The broader transseries volume supplies the formal-versus-analytic context
\cite{repo-volume}. Its own status inventory distinguishes theorems with
precise Lean counterparts from manuscript-level results. The present article
makes no status promotion for either that volume or the research package.
Only selected, relevant repository documents were inspected; this is not an
exhaustive audit of the repository.

\subsection{Relationship to standard methods}
The positive coefficient formula is classical Lagrange inversion, for which
Gessel's survey is a convenient reference \cite{gessel}. The idea of
condensation---one component carrying an unusually large share of a total---has
a substantial literature; Janson treats it for simply generated trees and
random allocations \cite{janson}. We do not claim to introduce condensation,
Lagrange inversion, Poisson approximation, or the saddle-point method.

The model here has the non-product weight
\((\sum_j\lambda_jm_j)^{\sum_jm_j-1}/\prod_jm_j!\).
Its large action is of order \(n/\log n\), not of order \(n\), and its
microscopic cloud has logarithmically growing size. The explicit equivalent,
the one-large-action/ones/twos reduction, and the resulting universality
constants are the model-specific refinements proved below. All necessary
estimates are supplied, rather than importing a random-allocation theorem
whose hypotheses have not been checked.

\section{Model, assumptions, and main theorems}
\label{sec:main}
Let \(\lambda_j\ge0\) be finite real numbers. We work first in
\(\R[[q]]\), with
\begin{equation}
 \U(q)=\sum_{j\ge1}q^j e^{\lambda_j\U(q)}
       =\sum_{n\ge1}u_nq^n.
 \label{eq:model}
\end{equation}
The identity is formal: no analytic convergence near \((q,u)=(0,0)\) is
presupposed. On writing \(q=e^{-x}\), it becomes a one-grid exponential
transseries \(\sum_nu_ne^{-nx}\). Well-based support and growth of its
coefficients are different questions.

The principal asymptotic hypothesis is
% ed. (2026-09-29): equation -> equation* for the \tag{Q} display, removing the
% duplicate hyperref destination equation.2.2; label, tag and numbering unchanged.
\begin{equation*}
 \boxed{\quad
 \lambda_j=j^2+\alpha j+\eta\frac{j}{\log j}
                  +o\!\left(\frac{j}{\log j}\right)
 \quad(j\longrightarrow\infty),\qquad \lambda_j\ge0.
 \quad}
 \tag{Q}\label{eq:Q}
\end{equation*}
The expression involving \(\log j\) is used only on a tail. The finitely many
initial slopes, especially \(\lambda_1\), are arbitrary nonnegative numbers.
Set
\begin{equation}
 b=1+\lambda_1,\qquad
 D(r)=2r^2+4r+1,
 \qquad r_n(r_n+1)e^{2r_n}=n,
 \label{eq:parameters}
\end{equation}
where the last equation has one positive real solution. Define
\begin{equation}
 x_n=\frac{n}{r_n+1},\qquad
 m_n=\frac{nr_n}{r_n+1},\qquad
 \sigma_n^2=\frac{m_n}{D(r_n)}.
 \label{eq:center}
\end{equation}
For later comparison,
\begin{align}
 r_n&=\frac12\log n-\log\log n+\log2
       +O\!\left(\frac{\log\log n}{\log n}\right),\label{eq:r-expansion}\\
 x_n&\sim\frac{2n}{\log n},\qquad
 \sigma_n^2\sim\frac{2n}{(\log n)^2}.
 \label{eq:rough-scales}
\end{align}
The implicit definition of \(r_n\), not a finite truncation of
\eqref{eq:r-expansion}, is to be used in the multiplicative formulas.

\begin{theorem}[Full coefficient equivalent]
\label[theorem]{thm:coefficient}
Under \eqref{eq:Q}, the formal solution exists uniquely, has \(u_1=1\) and
\(u_n\ge1\), and
\begin{equation}
 \boxed{\quad
 u_n\sim
 F_n:=\frac{\exp\!\left((2r_n+1)m_n+br_n^2
                              +\alpha r_n+\eta/2\right)}
                    {\sqrt{D(r_n)}}.
 \quad}
 \label{eq:main-equivalent}
\end{equation}
If \(\lambda_j=j^2\) for all sufficiently large \(j\), then
\(\alpha=\eta=0\) and the stronger bound
\begin{equation}
 u_n=F_n\left(1+O_\lambda\!\left(
                     \frac{(\log n)^2}{\sqrt n}\right)\right)
 \label{eq:rate}
\end{equation}
is valid. The constant may depend on the fixed exceptional slopes.
\end{theorem}

To state the probabilistic refinement, let
\begin{equation}
 \calM_n=\left\{\boldsymbol m=(m_1,m_2,\ldots):
                m_j\in\N,\ \sum_jjm_j=n\right\}.
 \label{eq:config-space}
\end{equation}
Here \(\N\) includes zero and only finitely many entries are nonzero. Put
\begin{equation}
 K=\sum_jm_j,\qquad S=\sum_j\lambda_jm_j,\qquad
 W(\boldsymbol m)=\frac{S^{K-1}}{\prod_jm_j!},\qquad
 \Prob_n(\boldsymbol m)=\frac{W(\boldsymbol m)}{u_n},
 \label{eq:probability}
\end{equation}
with \(0^0=1\). The coefficient formula below proves that this is a
probability distribution. Write
\(J=\max\{j:m_j>0\}\) and \(D_2=m_2\).

\begin{definition}
The discrete Gaussian \(G_n\) has law on \(\Z\)
\begin{equation}
 \Prob(G_n=k)=\frac{1}{Z_n}
       \exp\!\left(-\frac{(k-x_n)^2}{2\sigma_n^2}\right),\qquad
 Z_n=\sum_{k\in\Z}\exp\!\left(-\frac{(k-x_n)^2}{2\sigma_n^2}\right).
 \label{eq:discrete-gaussian}
\end{equation}
Let \(P_n\) be independent of \(G_n\), with distribution \(\Pois(r_n^2)\).
\end{definition}

\begin{theorem}[Microscopic condensation and joint local approximation]
\label[theorem]{thm:probability}
Under \eqref{eq:Q}, with probability
\begin{equation}
 1-O_\lambda\!\left(\frac{(\log n)^4}{n}\right),
 \label{eq:good-prob}
\end{equation}
there is exactly one part of size at least three. Its size is \(J\),
and all remaining parts have sizes one or two. Moreover,
\begin{equation}
 \TV\bigl(\mathcal L_n(J,D_2),\mathcal L(G_n,P_n)\bigr)\longrightarrow0.
 \label{eq:joint-TV}
\end{equation}
For eventually exact quadratic slopes the distance in \eqref{eq:joint-TV}
is \(O_\lambda((\log n)^2/\sqrt n)\).
Consequently,
\begin{equation}
 \left(\frac{J-x_n}{\sigma_n},\frac{D_2-r_n^2}{r_n}\right)
       \ \Longrightarrow\ (Z_1,Z_2),
 \label{eq:independent-normals}
\end{equation}
where \(Z_1,Z_2\) are independent standard normal random variables.
\end{theorem}

Thus a Poisson law, not merely a negligible relative error, describes the
small nonunit cloud. Its mean is asymptotic to \((\log n)^2/4\), while the
large-action variance is asymptotic to \(2n/(\log n)^2\).

\begin{corollary}[A three-level universality law]
\label[corollary]{cor:universality}
Let two nonnegative slope sequences satisfy \eqref{eq:Q}, with parameters
\((\lambda_1,\alpha,\eta)\) and
\((\widetilde\lambda_1,\widetilde\alpha,\widetilde\eta)\). Then
\begin{equation}
 \frac{u_n}{\widetilde u_n}\sim
 \exp\!\left((\lambda_1-\widetilde\lambda_1)r_n^2
       +(\alpha-\widetilde\alpha)r_n
       +\frac{\eta-\widetilde\eta}{2}\right).
 \label{eq:universality}
\end{equation}
In particular, finite changes to slopes with indices \(j\ge2\) do not alter
the multiplicative equivalent when the first slope is held fixed. A change
of the first slope does alter it, by the lognormal-size factor
\(e^{(\lambda_1-\widetilde\lambda_1)r_n^2}\).
\end{corollary}

\begin{theorem}[Refined positive-axis Borel growth]
\label[theorem]{thm:borel}
The integrated ordinary Borel transform
\begin{equation}
 B(t)=\sum_{n\ge1}\frac{u_n}{n!}t^n
 \label{eq:borel-def}
\end{equation}
is entire. Put \(s=s(t)=(\sqrt{1+4t}-1)/2\), so \(s(s+1)=t\). As
\(t\longrightarrow+\infty\),
\begin{equation}
 \boxed{\quad
 B(t)\sim
 \frac{\exp\!\left(s e^{2s}+bs^2+\alpha s+\eta/2\right)}
      {\sqrt{2s+1}}.
 \quad}
 \label{eq:borel-full}
\end{equation}
This is a multiplicative equivalent, not only an equivalent for a logarithm
or a double logarithm.
\end{theorem}

\begin{theorem}[Least-term equivalent]
\label[theorem]{thm:least}
For \(t>0\) let \(L(t)=\min_{n\ge1}u_nt^n\). As \(t\downarrow0\), put
\(R=\tfrac12\log(1/t)\). Then
\begin{equation}
 \boxed{\quad
 L(t)\sim
 \frac{\exp\!\left(-R^2/t+bR^2+\alpha R+\eta/2\right)}
      {\sqrt{2R^2+4R+1}}.
 \quad}
 \label{eq:least-full}
\end{equation}
Every minimizing index is asymptotic to
\begin{equation}
 N_0(t)=\frac{R(R+1)}{t}.
 \label{eq:least-index}
\end{equation}
Choosing a nearest integer to \(N_0(t)\) already gives a term asymptotic to
\(L(t)\).
\end{theorem}

\section{Formal solution and positive Lagrange weights}
\label{sec:lagrange}
For \(V\in q\R[[q]]\), every coefficient of
\(\sum_jq^je^{\lambda_jV}\) receives contributions only from \(j\le n\).
If \(V-W\) has order \(h\ge1\), then
\(e^{\lambda_jV}-e^{\lambda_jW}\) has order at least \(h\), and the
outer factor \(q^j\) raises it to at least \(h+1\). Picard iteration
therefore determines a unique formal solution coefficient by coefficient.
Its first coefficient is one.

\begin{lemma}[Coefficient identity]
\label[lemma]{lem:lagrange}
For every \(n\ge1\),
\begin{align}
 u_n&=\sum_{k=1}^n\frac1{k!}
       \sum_{\substack{j_1+\cdots+j_k=n\\j_i\ge1}}
                   (\lambda_{j_1}+\cdots+\lambda_{j_k})^{k-1}
       \label{eq:composition}\\
 &=\sum_{\boldsymbol m\in\calM_n}
       \frac{(\sum_j\lambda_jm_j)^{\sum_jm_j-1}}{\prod_jm_j!}.
       \label{eq:partition}
\end{align}
In particular all weights are nonnegative, and the \(k=1\) term gives
\(u_n\ge1\).
\end{lemma}
\begin{proof}
Introduce an auxiliary variable \(z\) and solve
\(V=z\Phi(q,V)\), where \(\Phi(q,v)=\sum_jq^je^{\lambda_jv}\).
Formal Lagrange inversion gives
\[
 V=\sum_{k\ge1}\frac{z^k}{k!}
      \left.\frac{\partial^{k-1}}{\partial v^{k-1}}
                   \Phi(q,v)^k\right|_{v=0}.
\]
At coefficient \(q^n\), only \(k\le n\) occurs, so setting \(z=1\) is
legitimate coefficientwise. Differentiating the product of exponentials
produces \eqref{eq:composition}. Grouping ordered compositions by their
multiplicities cancels the factor \(k!\) and yields \eqref{eq:partition}.
\end{proof}

For the unperturbed quadratic sequence the first coefficients are
\begin{equation}
 u_1=1,\quad u_2=2,\quad u_3=\frac{15}{2},\quad
 u_4=\frac{107}{3},\quad u_5=\frac{1555}{8},\quad
 u_6=\frac{17362}{15}.
 \label{eq:first-coefficients}
\end{equation}
In general \(u_2=1+\lambda_1=b\). Its reappearance as the coefficient of
\(r_n^2\) in the large-order equivalent is a concrete link between the first
nonlinear coefficient and the distant tail.

\section{Coarse localization, proved at exponential precision}
\label{sec:coarse}
This section supplies the localization needed for a self-contained proof.
It is a quadratic specialization of the envelope mechanism in the source
article, but no unproved concentration statement is needed subsequently.
Write \(\ell=\log n\) in estimates as \(n\to\infty\).

\begin{lemma}[Coarse large-action localization]
\label[lemma]{lem:coarse}
Suppose \(\lambda_j=j^2+O(j)\), with nonnegative slopes. With
\(R_n=n-K+1\), every fixed \(\epsilon>0\) satisfies
\begin{equation}
 \Prob_n\left(\left|\frac{R_n\ell}{n}-2\right|>\epsilon\right)
       \le e^{-c_\epsilon n+o(n)}.
 \label{eq:R-coarse}
\end{equation}
For every fixed \(\delta>0\),
\begin{equation}
 \Prob_n(J\le(1-\delta)R_n)\le e^{-c_\delta n+o(n)}.
 \label{eq:J-coarse}
\end{equation}
In particular, for any fixed \(1<a<2<c<\infty\),
\begin{equation}
 \Prob_n\left(J\notin\mathcal I_n\right)\le e^{-c_0n+o(n)},
 \qquad
 \mathcal I_n=\left[\frac{an}{\ell},\frac{cn}{\ell}\right]\cap\Z.
 \label{eq:macro-window}
\end{equation}
\end{lemma}
\begin{proof}
Choose \(C\) so that \(\lambda_j\le j^2+Cj\) for all \(j\).
For a composition with \(k\) parts, put \(R=n-k+1\). If
\(d_i=j_i-1\), then \(\sum_i d_i=R-1\) and
\[
 \sum_i j_i^2=\sum_i d_i^2+2(R-1)+k
              \le R^2+k-1.
\]
Thus the sum of slopes is at most \(R^2+C_1n\), where \(C_1\) is fixed.
The entire row with \(k=n-R+1\) parts is consequently bounded by
\begin{equation}
 T_n(R)=\binom{n-1}{R-1}
         \frac{(R^2+C_1n)^{n-R}}{(n-R+1)!}.
 \label{eq:row-envelope}
\end{equation}
A lower bound is obtained from one large part
\(J=\lfloor2n/\ell\rfloor\) and \(n-J\) ones:
\[
 u_n\ge\frac{\lambda_J^{n-J}}{(n-J)!}.
\]
Since \(\lambda_J=J^2(1+O(1/J))\), Stirling's formula gives
\begin{equation}
 \log u_n\ge n\ell-2n\log\ell+(\log4-1)n+o(n).
 \label{eq:coarse-lower}
\end{equation}

For completeness, the maximization of \eqref{eq:row-envelope} can be
localized without losing a constant of order \(n\). If \(R\ge\rho n\)
for a fixed \(\rho>0\), its logarithm is at most
\((1-\rho)n\ell+O(n\log\ell)+O(n)\), which is far below
\eqref{eq:coarse-lower}. A row within \(O(n)\) of that lower bound must
therefore have \(R/n\to0\). Its binomial entropy is then \(o(n)\).
The same comparison forces
\(\log(R^2+C_1n)\ge2\ell-O(\log\ell)\), so
\(R^2/n\to\infty\), \(\log R=\ell-O(\log\ell)\), and the additive
\(C_1n\) is negligible. A further comparison gives
\((R/n)\ell=O(\log\ell)\).

Put \(y=R\ell/n\). Uniform Stirling estimates in the just-localized rows
now yield
\begin{equation}
 \frac1n\log T_n(R)
 =\ell-2\log\ell+1+2\log y-y+o(1).
 \label{eq:coarse-functional}
\end{equation}
The function \(2\log y-y\) tends to minus infinity at either endpoint of
\((0,\infty)\) and has its unique maximum at \(y=2\).
On \(|y-2|\ge\epsilon\) it loses a fixed positive amount. The preliminary
localization shows that rows not covered by \eqref{eq:coarse-functional}
lose at least as much: otherwise a sequence of competing rows would have
satisfied those localization conditions. Summing at most \(n\) rows changes
a logarithm by at most \(\log n=o(n)\). This proves
\eqref{eq:R-coarse}, and also the matching upper version of
\eqref{eq:coarse-lower}.

For the largest part, always \(J\le R\). If \(J\le(1-\delta)R\), the
stronger estimate \(\sum_i d_i^2\le(J-1)(R-1)\) gives
\[
 \sum_i\lambda_{j_i}\le JR+C_2n.
\]
On the already localized rows, \(R^2\gg n\), so the right side is at most
\((1-\delta/2)R^2\) for large \(n\). Replacing the numerator of
\eqref{eq:row-envelope} by this capped value loses a factor
\((1-\delta/2)^{n-R}\), hence \(e^{-c n}\). The binomial entropy on those
rows is only \(e^{o(n)}\); it cannot cancel that loss. This proves
\eqref{eq:J-coarse}. Combining the two estimates, with slightly narrower
windows for \(R\), proves \eqref{eq:macro-window}.
\end{proof}

\section{The microscopic row estimate}
\label{sec:microscopic}
Fix \(J\in\mathcal I_n\), and write
\begin{equation}
 m=n-J,\qquad \Lambda=\lambda_J,\qquad
 t_J=\frac{m}{\Lambda},\qquad A_J=\frac{m^2}{\Lambda},\qquad
 C_J=\frac{\Lambda^m}{m!}.
 \label{eq:row-vars}
\end{equation}
In this section assume only \(\lambda_j=j^2+O(j)\). Uniformly on the
window, \(\Lambda\sim J^2\), \(t_J\asymp\ell^2/n\), and
\(A_J\asymp\ell^2\). Let \(u_{n,J}\) denote the sum of the weights with
largest part exactly \(J\).

\begin{lemma}[Uniform row equivalent and microscopic law]
\label[lemma]{lem:row}
Uniformly for \(J\in\mathcal I_n\),
\begin{equation}
 u_{n,J}=C_J e^{bA_J}
             \left(1+O_\lambda\!\left(\frac{\ell^4}{n}\right)\right).
 \label{eq:row-equivalent}
\end{equation}
Conditionally on the largest part being \(J\), the probability that there
is any additional part of size at least three is
\(O_\lambda(\ell^4/n)\). The conditional distribution of \(D_2\) is at
that same total-variation distance from \(\Pois(A_J)\).
\end{lemma}

\subsection{An upper bound that keeps the small-part entropy}
Distinguish one occurrence of the largest part \(J\). Let \(a_h\) be the
remaining multiplicities, so \(\sum_{h\le J}h a_h=m\). Write
\(r=\sum_h a_h\) and \(s=\sum_h\lambda_h a_h\). The weight with the
largest part marked is
\begin{equation}
 \frac{(\Lambda+s)^r}{\prod_{h\le J}a_h!}.
 \label{eq:marked}
\end{equation}
If there are several occurrences of \(J\), marking overcounts the original
weight by their number. Therefore a sum of marked weights is a valid upper
bound; when all remaining parts are ones or twos, it is exact.

Because \(s\ge0\) and \(r\le m\),
\[
 (\Lambda+s)^r\le\Lambda^r\exp(ms/\Lambda).
\]
Use
\[
 a_1=m-\sum_{h\ge2}h a_h,\qquad
 r=m-\sum_{h\ge2}(h-1)a_h,\qquad
 s=\lambda_1m+\sum_{h\ge2}(\lambda_h-\lambda_1h)a_h,
\]
and \(m!/a_1!\le m^{m-a_1}\). Dividing \eqref{eq:marked} by \(C_J\)
gives the upper bound
\begin{equation}
 e^{\lambda_1A_J}\prod_{h=2}^J\frac{\mu_h^{a_h}}{a_h!},
 \qquad
 \mu_h=\frac{m^h}{\Lambda^{h-1}}
           \exp\!\left(t_J(\lambda_h-\lambda_1h)\right).
 \label{eq:mu}
\end{equation}
We may drop the remaining mass constraint when summing a nonnegative upper
bound. Hence
\begin{equation}
 u_{n,J}\le C_J\exp\!\left(\lambda_1A_J+\sum_{h=2}^J\mu_h\right).
 \label{eq:upper-row}
\end{equation}

\begin{lemma}[Suppression of parts of size at least three]
\label[lemma]{lem:mu}
Uniformly on \(\mathcal I_n\),
\begin{equation}
 \mu_2=A_J+O_\lambda(\ell^4/n),\qquad
 \sum_{h=3}^J\mu_h=O_\lambda(\ell^4/n).
 \label{eq:mu-bounds}
\end{equation}
\end{lemma}
\begin{proof}
The first claim follows immediately from
\(\mu_2=A_Je^{t_J(\lambda_2-2\lambda_1)}\).
For the second choose a fixed \(C\) with
\(\lambda_h-\lambda_1h\le h^2+Ch\) for all \(h\), and bound \(\mu_h\)
by
\[
 v_h=\Lambda t_J^h e^{t_J(h^2+Ch)}.
\]
For \(3\le h\le J/2\),
\[
 \log(v_{h+1}/v_h)=\log t_J+t_J(2h+1+C)
 \le-(1-1/a+o(1))\ell.
\]
Thus that portion is geometrically decreasing and has sum
\(O(v_3)=O(\Lambda t_J^3)=O(\ell^4/n)\).
For \(J/2\le h\le J\), the logarithm of \(v_h\) is a convex quadratic
in \(h\). At the endpoints its leading terms are, respectively,
\(-J\ell/2+n/4+o(n)\) and \(-J\ell+n+o(n)\).
Since \(J\ell/n\ge a>1\), both are at most \(-c_1n\).
Convexity bounds the whole interval by the larger endpoint; its total
contribution is \(O(Je^{-c_1n})\). This proves \eqref{eq:mu-bounds}.
\end{proof}

Combining \cref{lem:mu} with \eqref{eq:upper-row} proves the upper half of
\eqref{eq:row-equivalent}. It also bounds the marked mass with at least one
additional part of size three or larger by
\begin{equation}
 C_J e^{\lambda_1A_J+\mu_2}
         \left(\exp\!\left(\sum_{h=3}^J\mu_h\right)-1\right)
 =C_Je^{bA_J}O_\lambda(\ell^4/n).
 \label{eq:bad-row}
\end{equation}
This bound includes configurations with a second largest part equal to \(J\).

\subsection{Matching lower bound and the Poisson cloud}
Retain exactly one \(J\), exactly \(d\) twos, and \(m-2d\) ones. Its
unmarked weight is
\begin{equation}
 w_{J,d}=\frac{\left(\Lambda+\lambda_1m
                          +(\lambda_2-2\lambda_1)d\right)^{m-d}}
                    {(m-2d)!\,d!},\qquad 0\le d\le m/2.
 \label{eq:exact-good-weight}
\end{equation}
For any fixed sufficiently large \(C_3\), uniformly for
\(0\le d\le C_3\ell^2\), expansion of the logarithm and of the falling
factorial gives
\begin{align}
 (m-d)\log\!\left(1+
        \frac{\lambda_1m+(\lambda_2-2\lambda_1)d}{\Lambda}\right)
   &=\lambda_1A_J+O_\lambda(\ell^4/n),\label{eq:log-expand}\\
 \frac{m!}{(m-2d)!\,m^{2d}}
   &=1+O(\ell^4/n).\label{eq:falling-expand}
\end{align}
For example, the error sizes in \eqref{eq:log-expand} are bounded by constants
times \(dm/\Lambda+m^3/\Lambda^2=O(\ell^4/n)\); the factorial error is
\(O(d^2/m)=O(\ell^4/n)\). Therefore
\begin{equation}
 w_{J,d}=C_Je^{\lambda_1A_J}\frac{A_J^d}{d!}
            \left(1+O_\lambda(\ell^4/n)\right).
 \label{eq:poisson-weights}
\end{equation}
Choose \(C_3\) larger than twice the largest asymptotic coefficient in
\(A_J/\ell^2\). The upper tail of \(\Pois(A_J)\) past \(C_3\ell^2\)
is \(e^{-c_2\ell^2}\), by the elementary exponential Markov bound
\(\Prob(P\ge k)\le\inf_{z>1}e^{A_J(z-1)}z^{-k}\).
That is smaller than \(\ell^4/n\). Summing \eqref{eq:poisson-weights}
proves the matching lower bound in \eqref{eq:row-equivalent}.

The same estimates, before summation, compare the normalized probabilities
on the cutoff with the Poisson probabilities in total variation. The cutoff
tail and \eqref{eq:bad-row} account for all remaining mass. This proves every
assertion of \cref{lem:row}.

\begin{remark}[Where the two contributions to \(b\) come from]
The factor \(e^{\lambda_1A_J}\) is the feedback of the many unit-action
factors. The factor \(e^{A_J}\) is their combinatorial replacement by twos.
Dropping the latter is not an innocuous simplification: it would lose a
factor of order \(e^{r_n^2}\). The numerical value of \(\lambda_2\) enters
only the smaller error at this level of precision.
\end{remark}

\section{The large-action saddle and the full equivalent}
\label{sec:saddle}
The microscopic estimate reduces the original sum over all integer
partitions to one positive sum over the largest action. We first analyze
the case of eventually exact quadratic slopes; then we show precisely how
the two additional tail parameters enter.

\subsection{An exact saddle coordinate}
For real \(0<x<n\), write
\begin{equation}
 H_n(x)=(n-x)\left[1+\log\frac{x^2}{n-x}\right],\qquad
 A_0(x)=\left(\frac{n-x}{x}\right)^2.
 \label{eq:H}
\end{equation}
Direct differentiation gives
\begin{align}
 H_n'(x)&=\log\frac{n-x}{x^2}+\frac{2(n-x)}x,\label{eq:Hprime}\\
 H_n''(x)&=-\frac1{n-x}-\frac4x-\frac{2(n-x)}{x^2}<0,
                      \label{eq:Hsecond}\\
 H_n'''(x)&=-\frac1{(n-x)^2}+\frac6{x^2}
                                      +\frac{4(n-x)}{x^3}.
                      \label{eq:Hthird}
\end{align}
There is exactly one critical point. Set \(r=(n-x)/x\) there.
Equation \eqref{eq:Hprime} becomes \(x=re^{2r}\), equivalently
\(n=r(r+1)e^{2r}\). Thus the critical point is exactly \(x_n\), and
\begin{equation}
 H_n(x_n)=(2r_n+1)m_n,\qquad
 -H_n''(x_n)=\frac{D(r_n)}{m_n}=\sigma_n^{-2}.
 \label{eq:saddle-identities}
\end{equation}
Taking logarithms in the defining equation for \(r_n\) first gives
\(r_n\sim\tfrac12\log n\). Substitution back into
\(2r_n=\log n-\log r_n-\log(r_n+1)\) gives
\eqref{eq:r-expansion} and then \eqref{eq:rough-scales}.

\subsection{A quantitative discrete Gaussian comparison}
For \(J\in\mathcal I_n\), Stirling's formula and \cref{lem:row} give,
when the slopes are eventually exactly quadratic,
\begin{equation}
 u_{n,J}=e^{H_n(J)}a_n(J)
               \left(1+O_\lambda(\ell^4/n)\right),\qquad
 a_n(x)=\frac{e^{bA_0(x)}}{\sqrt{2\pi(n-x)}}.
 \label{eq:row-saddle}
\end{equation}
Here and below the fixed macroscopic interval may be enlarged slightly to
contain each real segment between \(J\) and \(x_n\).
Uniformly on that interval,
\begin{equation}
 -H_n''(x)\asymp\frac{\ell^2}{n},\qquad
 |H_n'''(x)|=O(\ell^3/n^2),\qquad
 |(\log a_n)'(x)|=O_\lambda(\ell^3/n).
 \label{eq:saddle-controls}
\end{equation}
The last estimate follows from
\[
 (\log a_n)'(x)=-\frac{2bn(n-x)}{x^3}+\frac1{2(n-x)}.
\]
Let \(y=(J-x_n)/\sigma_n\) and
\(\epsn=\ell^2/\sqrt n\). Taylor's theorem, strict concavity, and
\eqref{eq:saddle-controls} imply
\begin{align}
 H_n(J)-H_n(x_n)&=-\tfrac12y^2+O(|y|^3/\sqrt n),
                                                        \label{eq:HTaylor}\\
 H_n(J)-H_n(x_n)&\le-c_1y^2,                 \label{eq:Hdomination}\\
 \left|\log\frac{a_n(J)}{a_n(x_n)}\right|
                       &\le C_1\epsn|y|.    \label{eq:amp-control}
\end{align}
These statements hold throughout the fixed macroscopic interval, not just
at a fixed value of \(y\).

To make the error after summation explicit, use
\(|e^v-e^w|\le |v-w|e^{\max(v,w)}\). Equations
\eqref{eq:HTaylor}--\eqref{eq:amp-control} then bound the difference of
the scaled summand and \(e^{-y^2/2}\) by
\begin{equation}
 C e^{-c_2y^2+C_2\epsn|y|}
          \left(\frac{|y|^3}{\sqrt n}+\epsn|y|\right).
 \label{eq:L1-saddle-error}
\end{equation}
The lattice spacing in \(y\) is \(\sigma_n^{-1}\to0\). Gaussian
moments, or a comparison on intervals of that length, show that the sum of
\eqref{eq:L1-saddle-error} is \(O_\lambda(\sigma_n\epsn)\).
The elementary Riemann-sum estimate
\begin{equation}
 \sum_{J\in\Z}e^{-(J-x_n)^2/(2\sigma_n^2)}
             =\sqrt{2\pi}\,\sigma_n+O(1)
 \label{eq:gaussian-sum}
\end{equation}
is uniform in the fractional part of \(x_n\). For instance, the error
between the integral and the sum of a continuously differentiable function
is bounded by the integral of the absolute derivative, which is bounded
for this Gaussian. The Gaussian mass outside \(\mathcal I_n\) is
exponentially small. Since \(1/\sigma_n=O(\ell/\sqrt n)=O(\epsn)\),
we obtain
\begin{equation}
 \sum_{J\in\mathcal I_n}u_{n,J}
 =e^{H_n(x_n)}a_n(x_n)\sqrt{2\pi}\,\sigma_n
                       \left(1+O_\lambda(\epsn)\right).
 \label{eq:sum-main}
\end{equation}
The row error \(O(\ell^4/n)=O(\epsn^2)\) is smaller. Finally,
\cref{lem:coarse} says that the left side is
\(u_n(1+O(e^{-c n+o(n)}))\); no absolute estimate of omitted rows is
being presumed. Substitution of \eqref{eq:saddle-identities} gives
\eqref{eq:main-equivalent} and \eqref{eq:rate} in the eventually
quadratic case.

\subsection{Linear and critical sublinear tail perturbations}
Assume now \eqref{eq:Q}, and let
\(\rho(x)=(n-x)/x\). Uniformly for \(J\in\mathcal I_n\),
\begin{align}
 (n-J)\log\frac{\lambda_J}{J^2}
   &=\alpha\rho(J)+\eta\frac{\rho(J)}{\log J}+o(1),
                                                   \label{eq:tail-expansion}\\
 \frac{(n-J)^2}{\lambda_J}-A_0(J)&=O_\lambda(\ell^3/n)=o(1).
                                                   \label{eq:A-perturb}
\end{align}
For the first equality, the quadratic Taylor error is
\(O((n-J)/J^2)=O(\ell^2/n)\). The unspecified tail error, when
multiplied by \((n-J)/J^2\), is
\(o((n-J)/(J\log J))=o(1)\) uniformly, because all the indices in the
window tend to infinity together and \(\rho(J)/\log J\) is bounded.
This uniformity is important: no regularity or differentiability of the
unspecified \(o\)-term is assumed.

Consequently \eqref{eq:row-saddle} holds with a relative \(1+o(1)\)
error and with the smooth amplitude
\begin{equation}
 a_n^*(x)=\frac{\exp\!\left(bA_0(x)+\alpha\rho(x)
                           +\eta\rho(x)/\log x\right)}
                   {\sqrt{2\pi(n-x)}}.
 \label{eq:perturbed-amplitude}
\end{equation}
The additional logarithmic derivatives are
\(O_\lambda(\ell^2/n)\), so the bound
\eqref{eq:amp-control} and the Gaussian comparison remain valid.
At the saddle,
\[
 \rho(x_n)=r_n,\qquad
 \log x_n=2r_n+\log r_n,\qquad
 \frac{r_n}{\log x_n}\longrightarrow\frac12.
\]
It follows that
\[
 a_n^*(x_n)=
 \frac{\exp(br_n^2+\alpha r_n+\eta/2)}{\sqrt{2\pi m_n}}
                       (1+o(1)).
\]
Repeating \eqref{eq:sum-main} proves \cref{thm:coefficient} in full.
For an arbitrary unspecified tail \(o(j/\log j)\), its modulus of
convergence contributes an unspecified \(o(1)\); the quantitative rate
\eqref{eq:rate} is therefore not asserted in that generality.

\begin{remark}[Why the saddle must be specified accurately]
An error small relative to \(x_n\) need not be small relative to
\(\sigma_n\). In particular, replacing \(x_n\) by
\(2n/\log n\) does not give a correctly centered Gaussian law. Similarly,
substituting a short logarithmic expansion for \(r_n\) into the large
exponent of \eqref{eq:main-equivalent} need not preserve a relative
\(1+o(1)\) error. The implicit scalar equation eliminates this avoidable
loss of precision.
\end{remark}

\section{The Gaussian--Poisson law and genuine local limits}
\label{sec:probability-proof}
The same estimates give distributional results without differentiating an
asymptotic equivalent. Normalize \eqref{eq:row-saddle} by the sum in
\eqref{eq:sum-main}. The summed absolute error
\eqref{eq:L1-saddle-error}, together with the relative error in the row
formula, proves
\begin{equation}
 \TV\bigl(\mathcal L_n(J),\mathcal L(G_n)\bigr)=o(1).
 \label{eq:J-TV}
\end{equation}
For eventually quadratic slopes the error is \(O_\lambda(\epsn)\).
The mass outside the macroscopic interval is already exponentially small.

\subsection{Removing the dependence of the Poisson mean on the large action}
Conditionally on \(J\in\mathcal I_n\), \cref{lem:row} compares
\(D_2\) with \(\Pois(A_J)\), with total-variation error
\(O_\lambda(\ell^4/n)\). For two nonnegative means, the elementary
Poisson coupling gives
\begin{equation}
 \TV(\Pois(v),\Pois(w))\le \min(1,|v-w|).
 \label{eq:poisson-coupling}
\end{equation}
Indeed, when \(v\le w\), add an independent \(\Pois(w-v)\) variable;
the resulting variables differ with probability at most \(w-v\).

For \(|J-x_n|=O(\sigma_n)\), equations
\eqref{eq:A-perturb} and
\(A_0'(x)=-2n(n-x)/x^3\) imply
\begin{equation}
 A_J-r_n^2=O_\lambda(\ell^2/\sqrt n)+o(1)=o(1).
 \label{eq:mean-freeze}
\end{equation}
In the eventually quadratic case, Gaussian domination from
\eqref{eq:Hdomination}--\eqref{eq:amp-control} gives the integrated bound
\[
 \E_n\!\left[\min(1,|A_J-r_n^2|);\ J\in\mathcal I_n\right]
                                  =O_\lambda(\epsn).
\]
In the general case the bounded quantity in brackets tends to zero in
probability, by \eqref{eq:J-TV} and \eqref{eq:mean-freeze}; hence its
expectation tends to zero. Mixture coupling, followed by
\eqref{eq:J-TV}, now proves \eqref{eq:joint-TV}, including its stated
rate in the eventually quadratic case.

The probability of an extra part of size at least three is
\(O_\lambda(\ell^4/n)\) on every macroscopic row by \cref{lem:row}.
The complement of those rows has exponentially small probability. This
proves \eqref{eq:good-prob}.
The standardized discrete Gaussian tends to a standard normal by
\eqref{eq:gaussian-sum}. The characteristic function of
\((P_n-r_n^2)/r_n\) is
\[
 \exp\!\left(r_n^2(e^{it/r_n}-1-it/r_n)\right)
                         \longrightarrow e^{-t^2/2}.
\]
Independence in the comparison law and vanishing total variation prove
\eqref{eq:independent-normals}. This completes the proof of
\cref{thm:probability}.

\subsection{Point probabilities in the central window}
Total-variation convergence alone is not a rescaled local limit theorem:
a small total error could in principle be concentrated at a few lattice
points. The available pointwise estimates rule that out here.

\begin{corollary}[Central two-dimensional local limit]
\label[corollary]{cor:local}
Under \eqref{eq:Q}, for every fixed \(M<\infty\), uniformly over integers
\(j,d\) with \(|j-x_n|\le M\sigma_n\) and
\(|d-r_n^2|\le Mr_n\),
\begin{equation}
 \sigma_nr_n\Prob_n(J=j,D_2=d)
 =\frac1{2\pi}\exp\!\left(-\frac{(j-x_n)^2}{2\sigma_n^2}
                         -\frac{(d-r_n^2)^2}{2r_n^2}\right)+o(1).
 \label{eq:local-limit}
\end{equation}
The same formula holds with the event replaced by
\(\{J=j,K=n+1-j-d\}\).
\end{corollary}
\begin{proof}
For configurations with only one part at least three,
\eqref{eq:poisson-weights} is a uniform relative estimate at every displayed
\(j,d\). The pointwise saddle expansion and the full coefficient equivalent
give the Gaussian factor. Stirling's formula for the Poisson probabilities,
with \(d=r_n^2+O(r_n)\), gives the second Gaussian factor uniformly.
Replacing \(A_j\) by \(r_n^2\) has vanishing relative effect there, by
\eqref{eq:mean-freeze}.

For a fixed central \(j\), all configurations with an additional part at
least three have probability at most
\(O_\lambda(\ell^4/n)\Prob_n(J=j)
=O_\lambda(\ell^4/(n\sigma_n))\).
After multiplication by \(\sigma_nr_n\), this is
\(O_\lambda(\ell^5/n)=o(1)\), even if all of that mass were at the
particular \(d\). Thus it cannot affect the limit.
On the good configurations, \(K=n+1-J-D_2\). The same bound handles
bad configurations for the second event as well.
\end{proof}

\begin{corollary}[The factor count and the excess action]
\label[corollary]{cor:K}
Let \(R_n=n-K+1\). Under \eqref{eq:Q},
\begin{equation}
 \frac{K-(n+1-x_n-r_n^2)}{\sigma_n}\Longrightarrow -Z_1,
 \qquad
 \frac{R_n-(x_n+r_n^2)}{\sigma_n}\Longrightarrow Z_1.
 \label{eq:K-limit}
\end{equation}
Moreover, \(R_n-J\) is at total-variation distance \(o(1)\) from
\(\Pois(r_n^2)\).
\end{corollary}
\begin{proof}
On the event of \eqref{eq:good-prob}, \(R_n=J+D_2\) exactly. Apply the
joint approximation and use \(r_n/\sigma_n\to0\).
\end{proof}
The correction \(r_n^2\) is invisible on the larger Gaussian scale but
is indispensable for the microscopic difference \(R_n-J\). This explains
why the earlier statement \(J/R_n\to1\) contains much less information.

\section{Universality, perturbations, and what the constants mean}
\label{sec:universality}
Dividing the equivalents in \cref{thm:coefficient} proves
\cref{cor:universality}. It is useful to separate the three mechanisms.

\subsection{Finite-prefix universality and its exception}
Suppose the slopes are eventually exactly quadratic. Every fixed
\(\lambda_h\), \(h\ge2\), influences lower-order terms in the
microscopic row expansion, but none influences its leading equivalent.
This is not because small parts are absent: the number of twos tends to
infinity in probability. Rather, the large slope \(\lambda_J\asymp J^2\)
dominates the feedback sum so strongly that the numerical contribution
of each fixed nonunit slope is negligible in its exponent.

The first slope is different. There are about \(n\) action-one factors,
so their cumulative feedback contributes
\(\lambda_1 n^2/J^2\sim\lambda_1r_n^2\).
For the baseline \(\lambda_j=j^2\), one has \(b=2\). Keeping that tail
but setting \(\lambda_1=c\ge0\) gives
\begin{equation}
 \frac{u_n^{(c)}}{u_n^{(1)}}\sim e^{(c-1)r_n^2}.
 \label{eq:first-slope-example}
\end{equation}
Thus equality of the tail of the defining slope sequence does not imply
equality of coefficient equivalents unless the first slope is also fixed.

\subsection{The linear tail is visible at a smaller exponential scale}
For \(\lambda_j=j^2+\alpha j+o(j/\log j)\), the extra factor is
\(e^{\alpha r_n}\). The exact relation defining \(r_n\) rewrites it as
\begin{equation}
 e^{\alpha r_n}
       =\left(\frac{n}{r_n(r_n+1)}\right)^{\alpha/2}.
 \label{eq:linear-tail-factor}
\end{equation}
It is therefore polynomial in \(n\), with a logarithmic correction.
There is no need to re-center the leading Gaussian by an amount comparable
to its standard deviation: the induced saddle shift is negligible on that
scale, although the amplitude itself is not negligible.

\subsection{A sharp constant-scale perturbation}
A perturbation \(\eta j/\log j\) produces the finite factor
\(e^{\eta/2}\), with the first slope and the linear coefficient held
fixed. The mechanism is the exact evaluation
\[
 \frac{n-J}{J\log J}\sim\frac{r_n}{\log x_n}\longrightarrow\frac12.
\]
Any tail \(o(j/\log j)\) disappears from the multiplicative equivalent.
This gives an actual stability theorem, not a formal Taylor heuristic.
For instance, changing the tail from \(j^2\) to
\(j^2+2j/\log j\), while preserving \(\lambda_1\), multiplies the
large-order coefficients asymptotically by \(e\).

\begin{remark}[Limits of this universality statement]
The hypotheses deliberately specify the tail at the scale \(j/\log j\).
The weaker assumption \(\lambda_j\sim j^2\) is not sufficient for the
same equivalent: a perturbation such as \(j^2/\sqrt{\log j}\) is
relatively small but is multiplied by a large exponent in the coefficient
formula. Nor does this theorem cover sign cancellations or negative slopes.
\end{remark}

\section{A second saddle: the entire Borel transform}
\label{sec:borel-proof}
\subsection{Entirety and the continuous coefficient envelope}
The coefficient equivalent gives
\[
 \left(\frac{u_n}{n!}\right)^{1/n}
                         \sim\frac4{(\log n)^2}\longrightarrow0.
\]
The root test proves that \(B\) in \eqref{eq:borel-def} is entire.
For real \(v>0\), define \(r=r(v)\) by
\(v=r(r+1)e^{2r}\), and let
\begin{equation}
 F_0(v)=(2r+1)\frac{vr}{r+1}.
 \label{eq:F0}
\end{equation}
Implicit differentiation gives the exact identities
\begin{equation}
 r'(v)=\frac{r(r+1)}{vD(r)},\qquad F_0'(v)=2r.
 \label{eq:r-derivative}
\end{equation}
Applying Stirling's formula to \(n!\), the positive Borel summand is
asymptotic, uniformly as the integer \(n\to\infty\), to
\begin{equation}
 e^{K(n;t)}c(n),\qquad
 c(v)=\frac{\exp(br(v)^2+\alpha r(v)+\eta/2)}
                 {\sqrt{2\pi vD(r(v))}},
 \label{eq:borel-summand}
\end{equation}
where
\begin{align}
 K(v;t)&=F_0(v)-v\log v+v+v\log t\notag\\
       &=v\left[\frac1{r(v)+1}
                      +\log\frac{t}{r(v)(r(v)+1)}\right].
 \label{eq:K}
\end{align}
The meaning of uniformity here is simple: the error in the coefficient
equivalent is a sequence tending to zero, so its supremum over all
\(n\ge N\) tends to zero. Positivity permits this error to be passed
through any tail sum.

\subsection{Exact location and curvature of the Borel saddle}
Equations \eqref{eq:r-derivative} and \eqref{eq:K} imply
\begin{equation}
 K'(v;t)=\log\frac{t}{r(v)(r(v)+1)},\qquad
 K''(v;t)=-\frac{2r(v)+1}{vD(r(v))}<0.
 \label{eq:Kderivatives}
\end{equation}
Let \(s(s+1)=t\). The unique critical point, its value, and its variance
parameter are therefore
\begin{equation}
 v_* =t e^{2s},\qquad
 K(v_*;t)=s e^{2s},\qquad
 \nu^2=\frac{v_*D(s)}{2s+1}.
 \label{eq:borel-saddle}
\end{equation}
In particular \(\nu/v_*\to0\). Uniformly for
\(v\in[v_*/2,2v_*]\), one has \(r(v)=s+O(1)\), and direct
differentiation gives
\begin{equation}
 |K'''(v;t)|=O(1/(v_*^2s)),\qquad
 |(\log c)'(v)|=O_\lambda(s/v_*).
 \label{eq:second-saddle-controls}
\end{equation}
After multiplication by \(\nu^3\) and \(\nu\), respectively, the
right sides become \(O(\sqrt{s/v_*})\) and
\(O_\lambda(s^{3/2}/\sqrt{v_*})\), both tending to zero.
Consequently Taylor's formula and a discrete Gaussian comparison on this
interval yield
\begin{equation}
 \sum_{v_*/2\le n\le2v_*} e^{K(n;t)}c(n)
           \sim e^{K(v_*;t)}c(v_*)\sqrt{2\pi}\,\nu.
 \label{eq:borel-central-sum}
\end{equation}
This is the same estimate as in \cref{sec:saddle}, now with
\(\nu\to\infty\) and an even more slowly varying relative amplitude.

\subsection{Tail control for the infinite positive sum}
We record the tail argument because a saddle estimate on a finite window
alone does not establish an entire-function growth theorem.
On \([v_*/2,2v_*]\), strict concavity loses at least
\(c v_*/s\) at either endpoint. Also
\(K'(v_*/2;t)\ge c/s\) and \(K'(2v_*;t)\le-c/s\), for large \(s\),
by \eqref{eq:Kderivatives} or by integrating its second derivative.
For \(n\le v_*/2\), the core exponent is bounded by its endpoint value,
since it increases up to \(v_*\). The amplitude has logarithm bounded
above by \(O_\lambda(1+s^2)\) on \(1\le n\le v_*/2\).
Even after multiplication by the number of terms, this tail is negligible
compared with the central sum.

For \(v\ge2v_*\), concavity gives
\[
 K(v;t)\le K(2v_*;t)-\frac{c}{s}(v-2v_*).
\]
From \eqref{eq:r-derivative},
\(|(\log c)'(v)|\le C_\lambda(1+\log v)/v\) for large \(v\).
The supremum of this bound over \(v\ge2v_*\) is
\(O_\lambda(s/v_*)=o(1/s)\). Thus half of the displayed linear
decay absorbs all possible amplitude growth. The remaining tail is a
geometric sum with factor of order \(s\), multiplied by an endpoint
already smaller by \(e^{-c v_*/s}\).
Finitely many initial Borel terms grow only polynomially in \(t\), and
are negligible. These bounds also justify insertion of the uniform
coefficient equivalent into the infinite tail sum.

Substitution of \eqref{eq:borel-saddle} into
\eqref{eq:borel-central-sum} now gives
\[
 e^{s e^{2s}}
 \frac{e^{bs^2+\alpha s+\eta/2}}{\sqrt{2\pi v_*D(s)}}
 \sqrt{\frac{2\pi v_*D(s)}{2s+1}},
\]
which is exactly \eqref{eq:borel-full}. This proves \cref{thm:borel}.

\subsection{Consequences and directional limitations}
Since \(s=\sqrt t-\tfrac12+O(t^{-1/2})\),
\begin{equation}
 \log B(t)\sim s e^{2s}
             \sim\frac{\sqrt t}{e}e^{2\sqrt t},\qquad
 \log\log B(t)\sim2\sqrt t.
 \label{eq:borel-logconsequences}
\end{equation}
The first equivalent and the full formula retain information lost by the
second. With fixed tail parameters, changing \(\lambda_1\) by
\(\Delta\) changes \(B(t)\) by a multiplicative factor asymptotic to
\(e^{\Delta s^2}\), although it leaves the leading equivalent for
\(\log B\) unchanged.

For every \(q>0\),
\begin{equation}
 \frac1q\int_0^\infty e^{-t/q}B(t)\dd t=+\infty.
 \label{eq:laplace-failure}
\end{equation}
Indeed, \eqref{eq:borel-logconsequences} gives
\(\log B(t)-t/q\to+\infty\); the positive integrand does not even
tend to zero. An entire Borel transform does not automatically have the
growth bounds needed for directional Borel--Laplace summation; this
distinction is part of the standard summability framework \cite{sauzin}.
Nothing in this positive-axis proof determines growth on the negative axis,
angular cancellation, or a Stokes constant.

There is also a direct obstruction to a positive-real literal solution of
\eqref{eq:model}. If \(q\in(0,1)\) and \(u>0\), then
\(j\log q+\lambda_j u\to+\infty\), so
\(q^je^{\lambda_j u}\) fails to tend to zero. Any proposed accelerated
positive-direction value must therefore specify a continuation or
regularization of the kernel; it cannot silently be called a solution of
that divergent literal sum.

\section{Least-term asymptotics without an unwarranted remainder claim}
\label{sec:least-proof}
For each fixed \(t>0\), \cref{thm:coefficient} implies
\(u_n^{1/n}\to\infty\). Thus \(u_nt^n\to\infty\), and its minimum
is attained at a finite index.
Put \(R=\tfrac12\log(1/t)\) and use the continuous envelope
\begin{equation}
 Q(v;t)=F_0(v)+v\log t+g(r(v)),\qquad
 g(r)=br^2+\alpha r+\eta/2-\tfrac12\log D(r).
 \label{eq:least-envelope}
\end{equation}
By \eqref{eq:r-derivative}, the derivative of the core
\(F_0(v)+v\log t\) is \(2r(v)-2R\). Its unique minimum is at
\(N_0=R(R+1)e^{2R}=R(R+1)/t\), where the core value is exactly
\(-R^2/t\).

On a fixed relative neighborhood of \(N_0\), the core second derivative
is \(2r'(v)\asymp1/N_0\), while
\[
 \frac{\partial}{\partial v}g(r(v))=O_\lambda(R/N_0),\qquad
 \frac{\partial^2}{\partial v^2}g(r(v))=O_\lambda(R/N_0^2).
\]
It follows by the mean value theorem that the minimum of the smooth
full envelope is displaced from \(N_0\) by \(O_\lambda(R)\).
The resulting change in its minimum value is
\(O_\lambda(R^2/N_0)=O_\lambda(t)=o(1)\).
Rounding either this minimum or \(N_0\) to an integer changes the value
by \(o(1)\).

For clarity, the local minimum cannot be bypassed at a far-away index.
The full envelope is strictly convex for all sufficiently large \(v\):
its positive core curvature is asymptotic to \(1/v\), and the amplitude
curvature is \(O_\lambda((\log v)/v^2)\). Its derivative has opposite
signs at \(N_0/2\) and \(2N_0\) when \(t\) is small. Therefore it has
one eventual minimum. On the complement of any fixed relative
neighborhood of \(N_0\), integration of the core curvature gives a
loss of order \(N_0\), whereas the amplitude variation is only
\(O_\lambda(R^2)\) on fixed relative windows; convexity continues that
loss beyond them. The finitely many smaller indices give terms of
polynomial size in \(t\), much larger than \(e^{-R^2/t+O(R^2)}\).

The coefficient equivalent means that
\(\log(u_nt^n)=Q(n;t)+o(1)\), uniformly over all sufficiently large
integers. Thus it may be minimized without losing more than \(o(1)\)
in the minimum value. Every minimizing index tends to infinity and must
satisfy \(n/N_0\to1\), by the preceding convexity loss.
Evaluating \(g\) at \(r=R\) and exponentiating proves
\cref{thm:least}, including its assertion about a nearest integer to
\(N_0\).

\begin{remark}[A scale, not yet an analytic error theorem]
The dominant least-term scale is
\begin{equation}
 \exp\!\left(-\frac{\log^2(1/t)}{4t}\right).
 \label{eq:superleast}
\end{equation}
It is smaller than \(e^{-A/t}\) for every fixed \(A>0\), as
\(t\downarrow0\). This does not imply that a sectorial solution has a
remainder of this size. A remainder theorem must relate the actual function
to its truncated Taylor series and control cancellations and analytic
continuation. The source article provides a left-sector fixed point, but
its Poincar\'e statement alone supplies no such uniform optimal-order
estimate \cite{repo-classification}.
\end{remark}

\section{Exact computation and reproducibility}
\label{sec:computation}
The accompanying program is designed to audit finite algebra and to show
how slowly the asymptotic regime can be approached. Its exact calculations
are logically separate from the asymptotic proofs.

\subsection{An integral coefficient recurrence}
For integer slopes, let
\begin{equation}
 b_n=n!u_n,\qquad E_{j,m}=m![q^m]e^{\lambda_j\U(q)}.
 \label{eq:scaled-recurrence-defs}
\end{equation}
Differentiation of the exponential identity gives
\begin{align}
 E_{j,0}&=1,\notag\\
 E_{j,m}&=\lambda_j\sum_{k=1}^m
            \binom{m-1}{k-1}b_kE_{j,m-k},\qquad m\ge1,
                                                   \label{eq:exact-rec-E}\\
 b_n&=\sum_{j=1}^n\frac{n!}{(n-j)!}E_{j,n-j}.
                                                   \label{eq:exact-rec-b}
\end{align}
When computing \(b_n\), the required exponential coefficients have
\(m\le n-1\), so only previously computed \(b_k\) are used. Starting
from \(E_{j,0}=1\), every operation in these formulas is an integer
operation. This gives an independent implementation of the formal fixed
point without enumerating partitions.

A second implementation enumerates integer partitions of \(n\) and sums
\eqref{eq:partition} using exact rational arithmetic. Agreement of these
two implementations is a useful test because their combinatorial
organization and normalization differ substantially.

\subsection{Executed checks}
The program computed exact coefficients through degree \(256\) for four
models: the unmodified quadratic slopes, and the three modifications
\(\lambda_1=0\), \(\lambda_1=3\), and \(\lambda_2=7\), with all other
slopes unchanged. For each model, all degrees \(1\) through \(18\) were
also computed by independent partition enumeration. All \(72\) exact
comparisons passed.

For the quadratic model it additionally summed the exact weights
\eqref{eq:exact-good-weight} over all possible \(J\ge3\) and \(d\).
The resulting mass was checked to lie between zero and the full
coefficient. The reported probabilities are evaluations of ratios of
these exact integers. The asymptotic formula, logarithms, and displayed
ratios were evaluated at \(80\)-decimal working precision using
\texttt{mpmath}; they are not outward-rounded interval bounds.

\begin{table}[htbp]
\centering
\caption{Quadratic slopes: the exact coefficient divided by the leading
 equivalent, and the probability of one large action with only ones and twos.}
\label{tab:diagnostics}
\begin{tabular}{@{}rrrr@{}}
\toprule
\(n\)&\(r_n\)&\(u_n/F_n\)&Good-configuration probability\\
\midrule
20  &1.087760&3.531186&0.561985\\
40  &1.298005&2.090003&0.767695\\
80  &1.519691&1.689076&0.849825\\
120 &1.654076&1.557336&0.876061\\
180 &1.791611&1.455860&0.896387\\
256 &1.913461&1.384043&0.910996\\
\bottomrule
\end{tabular}
\end{table}
At degree \(256\), the exact coefficient is still about \(38\%\) larger
than the leading equivalent \(F_n\). The theorem
is an asymptotic statement, and its conservative
\(O((\log n)^2/\sqrt n)\) rate is not a small numerical bound at those
orders. The observed values support neither a monotonic-convergence claim
nor an explicit finite-order error constant. Those would require separate
arguments.

\begin{table}[htbp]
\centering
\caption{Finite-perturbation ratios. Each displayed column tends to one
 by \cref{cor:universality}; the finite values are diagnostics, not proofs.}
\label{tab:perturbations}
\begin{tabular}{@{}rrrr@{}}
\toprule
\(n\)&\(u_n^{[\lambda_2=7]}/u_n\)&
\(e^{-2r_n^2}u_n^{[\lambda_1=3]}/u_n\)&
\(e^{r_n^2}u_n^{[\lambda_1=0]}/u_n\)\\
\midrule
20  &3.642423&6.041575&0.888086\\
40  &1.897724&2.674726&0.960560\\
80  &1.538280&1.699454&0.996221\\
120 &1.436076&1.516262&1.005374\\
180 &1.358496&1.394877&1.010625\\
256 &1.303662&1.317723&1.013079\\
\bottomrule
\end{tabular}
\end{table}

\subsection{Files and rebuild procedure}
The package contains the article source and PDF, the verification program,
exact coefficient tables, numerical diagnostics, a machine-readable check
report, source-provenance notes, and a build script. From the package root,
run
\begin{quote}\small\ttfamily
python -m pip install -r requirements.txt\\
python code/verify.py --order 256 --check-order 18\\
sh build.sh
\end{quote}
The exact recurrence and partition calculations use Python's standard
library. The listed \texttt{mpmath} dependency is used for diagnostics.
The PDF needs a standard LaTeX installation with the packages named in the
source. No internet retrieval, external figure, or external bibliography
file is required to rebuild it.

\section{Proof status and a feasible formalization route}
\label{sec:formalization}
Every main assertion above has a conventional mathematical proof in this
article. None has been machine-checked in Lean in this work. The finite
program checks do not certify infinite tails, limiting distributions, or
Borel growth.

A useful formalization can be staged without implementing a complete
transseries field. First formalize the coefficientwise contraction and the
finite coefficient identity in a formal-power-series ring. The next stage
is finite: multiplicity vectors, the marked-largest-part inequality, the
exact ones/twos weights, and the positive exponential majorant. These are
natural targets for algebraic and combinatorial proof automation.

The central analytic interface should be the uniform row lemma, with its
parameters and macroscopic window made explicit. A reusable discrete
Laplace theorem should accept bounds on the second and third derivatives,
a bound on the logarithmic derivative of the amplitude, and a certified
outside-window estimate. The proof in \cref{sec:saddle} identifies exactly
those hypotheses. Such an interface would also support the second Borel
saddle, rather than requiring a special-purpose proof for each sum.

Finally, probability statements should be formulated in total variation
before passing to weak limits. The Poisson coupling and the uniform
pointwise bounds are simpler, more informative intermediate goals than an
isolated central limit theorem. The distinction between coefficient
asymptotics and analytic remainders should remain explicit in the API;
there is no valid implication from \cref{thm:least} alone to a remainder
certificate.

\begin{center}
\small
\begin{tabular}{@{}p{0.27\textwidth}p{0.61\textwidth}@{}}
\toprule
Result & Dependencies and status\\
\midrule
Positive coefficient identity & Classical Lagrange inversion, with a finite
coefficientwise specialization; rederived here and checked at finite orders.\\
Coarse localization & Positive envelope, Stirling estimates, one-dimensional
strict concavity; rederived from the source article's mechanism.\\
Microscopic row lemma & Marking, exponential majorization, a small-part tail
bound, and exact ones/twos weights; proved here.\\
Full coefficient equivalent & Microscopic row lemma, explicit saddle,
Gaussian sum, and perturbation expansion; proved here.\\
Joint and local limits & Pointwise row weights, summed absolute saddle error,
Poisson coupling, and Stirling; proved here.\\
Borel equivalent & Coefficient equivalent, a second saddle, and control of
both infinite tails; proved here.\\
Least term & Coefficient equivalent, convex minimization, integer rounding;
proved here, but not an analytic remainder theorem.\\
\bottomrule
\end{tabular}
\end{center}

\section{Further research questions}
\label{sec:questions}
The following questions continue the proved results rather than restating
questions already settled in this article. Assertions labeled as heuristics
are not theorems of this paper.

\subsection{All-order corrections and the first appearance of each small slope}
For eventually quadratic slopes, determine a systematic expansion of
\(u_n/F_n\) in inverse powers of \(n\), with coefficients expressed in
\(r_n\) and the fixed initial slopes. The exact microscopic formula
provides a route: expand the falling factorial, the feedback logarithm,
and the discrete saddle to matching orders. The first omitted primitive
cloud has scale \(m^3/\Lambda^2\asymp(\log n)^4/n\).
At what precise order does each \(\lambda_h\), \(h\ge2\), first
become visible, and which coefficients admit sign or monotonicity
statements? A useful answer should include a uniform truncation bound, not
only a list of formal correction polynomials.

\subsection{The cloud hierarchy for nonquadratic slopes}
For \(\lambda_j\sim a j^p\), \(p>1\), the source article proves only
macroscopic concentration and logarithmic coefficient growth. At a
one-large-action configuration, the natural formal small-part activity is
approximately
\[
 \mu_h\approx\frac{n^h}{(aJ^p)^{h-1}},\qquad
 J\asymp n/\log n.
\]
Thus its power of \(n\) is \(p-(p-1)h\). This is a heuristic predicting
changes at \(p=h/(h-1)\): for \(p>2\), even the twos should disappear;
for \(1<p<2\), several small sizes may occur in large numbers. At a
critical exponent the logarithmic factors still matter. Prove the
appropriate multi-cloud limit theorem, including interactions and the
correct shifted saddle. The present quadratic theorem proves one member
of that hierarchy, not the hierarchy itself.
% ed. (2026-09-29): partly answered by the finite-core and Poisson-layer packages.
\begin{ednote}
Partly answered for $\lambda_j=aj^p$ (after finitely many nonnegative
changes) by two other packages of this series.
\nolinkurl{Finite_Core_Universality_Exponential_Feedback/} proves that
an analytic core of $M$ actions gives coefficients asymptotic to $u_n$
exactly when $M(p-1)>1$ (its Theorem~\texttt{thm:cutoff}), so for $p>2$ a
one-action core suffices, as the heuristic above predicts, with a
multiplicative saddle equivalent (its Theorem~\texttt{thm:saddle}).
\nolinkurl{Poisson_Layers_Finite_Core_Boundary/} proves Poisson and
Gaussian laws for the number of occurrences of the first omitted action,
jointly with the largest action, throughout $(M+1)(p-1)>1$ (its
Theorems~\texttt{thm:probability} and \texttt{thm:joint}). The joint law of the whole cloud
vector is not treated.
\end{ednote}

\subsection{Boundary perturbations without a smooth asymptotic tail}
The scale \(j/\log j\) is now known to affect the constant in the
coefficient equivalent. Classify tails of the form
\(\lambda_j=j^2+\alpha j+(j/\log j)\eta_j\), where \(\eta_j\)
oscillates instead of converging. Slow oscillations suggest an amplitude
sampled near \(j=x_n\); oscillations across a window of width
\(\sigma_n\) may instead be averaged by a Gaussian profile. Determine
necessary and sufficient local regularity hypotheses, and construct
examples with genuinely nonconvergent normalized coefficients.

\subsection{Compositional inverses and cancellation-sensitive large order}
Let \(\widehat Q=\U^{\langle-1\rangle}\). The source article transfers
Gevrey-class membership but does not determine exact type or a coefficient
equivalent. Does the one-large-action mechanism survive the signed
reversion weights after cancellation? Determine the first exponential and
subexponential scales of the inverse coefficients, and identify when the
first-slope factor remains visible. Positivity of \eqref{eq:partition}
cannot simply be reused for a signed inversion formula.
% ed. (2026-09-29): exact type answered by the weighted-type package; coefficient equivalent claimed by two batch-49 packages.
\begin{ednote}
The exact type is answered by
\nolinkurl{Sharp_Weighted_Type_Formal_Reversion/}: reversion preserves the
Gevrey type (its Theorem~\texttt{thm:inversefeedback}) and the refined inverse type
is $4$ for $\lambda_j=j^2$ (its Corollary~\texttt{cor:quadratic}). A coefficient
equivalent for $\lambda_j=aj^2$ is claimed, independently and not yet
reviewed, by \nolinkurl{Quadratic_Exponential_Feedback_After_Reversion/}
and \nolinkurl{Signed_Quadratic_Feedback_Inversion/} (Theorem~\texttt{thm:main} in
each): $v_n\sim-S_n\exp(-(1+1/a)r_n)$ in the notation of the finite-core
package; the first of them also treats slope tails $aj^2+dj+e$ and
$\kappa j/\log j$ perturbations, and in both the first-slope factor
enters through $\lambda_1+w_2$.
\end{ednote}

\subsection{Angular growth and negative-direction Borel summation}
The positive-axis Borel equivalent is now explicit, but it gives no
cancellation estimate for \(B(-t)\) or for nearby complex rays. Determine
sectorial growth, possible indicator transitions, and whether a suitable
negative-direction Borel--Laplace sum agrees with the literal left-sector
fixed point. The proof must establish the analytic identification as well
as growth; an entire Borel transform alone settles neither.
% ed. (2026-09-29): answered in the fine sense by the negative-ray package and negatively in the angular sense by the natural-boundaries package.
\begin{ednote}
Answered in two senses by later packages.
\nolinkurl{Negative_Ray_Summation_Exponential_Feedback/} (its
Theorem~\texttt{thm:threshold}) proves, for unit amplitudes and
$\lambda_j=O((j\log j)^2)$, which contains \eqref{eq:Q}, fine
(fixed-width half-strip) Borel--Laplace summability of $\U$ and of its
inverse in direction $\pi$, with sums that satisfy the literal kernel near
the negative axis.
\nolinkurl{Natural_Boundaries_Quadratic_Exponential_Feedback/} proves that
for $\lambda_j=j^2$ (its Theorem~\texttt{thm:main}), and for slopes that are
eventually $j^2+\beta j+\gamma$ (its Theorem~\texttt{thm:rational-main}), no open
arc of directions around $\pi$ gives angular $1$-summability, because the
literal inverse has a natural boundary on the imaginary diameter.
\end{ednote}

\subsection{Acceleration adapted to the double-exponential growth}
Construct and analyze an acceleration procedure adapted to
\(\log B(t)\sim(\sqrt t/e)e^{2\sqrt t}\). Specify the analytic
continuation or regularization of the countable kernel that the proposed
sum solves. Different choices may be inequivalent, since the literal
positive-real kernel diverges for every positive candidate value. A
complete result should identify uniqueness conditions and compatibility
with composition or inversion.

\subsection{From the least term to the actual sectorial remainder}
The proved least-term scale \eqref{eq:superleast} suggests a refined
optimal-order problem for the left-sector analytic solution. Does its
Taylor remainder at order near \(R(R+1)/t\) have that scale, a different
scale, or additional oscillatory cancellation? Establish a uniform
remainder theorem on proper sectors, and only then compare it with the
least term. This is a sharper question than coefficient Gevrey membership.
% ed. (2026-09-29): uniform remainder part answered by the sectorial-summability package; comparison with the least term open.
\begin{ednote}
The uniform remainder theorem is supplied by
\nolinkurl{Sharp_Negative_Direction_Summability_Exponential_Feedback/}:
strong Gevrey-$1$ remainders on every proper left sector (its
Theorem~\texttt{thm:headline}) and an actual truncation error at most
$C\exp[-c\log^2(1/|q|)/|q|]$ at an index of order $\log^2(1/|q|)/|q|$
(its Corollary~\texttt{cor:actualerror}), with an unspecified constant $c$; the
negative-ray package bounds the inverse on the same scale (its
Theorem~\texttt{thm:optimal}, an upper bound). Whether the actual remainder has the
scale of the least term remains open.
\end{ednote}

\subsection{Amplitudes, noninteger actions, and several feedback components}
Replace the model by
\[
 U=\sum_j c_jq^{a_j}e^{\lambda_jU},\qquad c_j\ge0,
\]
first with integer actions and then with well-based real actions. Which
amplitude and action assumptions preserve a one-large-action reduction?
The unit-action reservoir is essential to the present balance. For vector
feedback, determine whether an analogous reduction is governed by a
nonnegative matrix or by several competing large components. Signed or
complex amplitudes require a separate noncancellation theory.

\subsection{Explicit finite certificates and scalable computation}
The current relative error constants are asymptotic and depend on the
fixed exceptional slopes. Extract computable thresholds and constants
from the microscopic and saddle bounds, use directed interval arithmetic,
and give finite enclosures for \(u_n\) at orders beyond exact recurrence
computation. A particularly useful tool would certify both the coefficient
and the probability of exceptional configurations, with its proof
obligations structured for eventual Lean verification.

\section{Conclusion}
The unresolved finer-asymptotics question in the repository's
exponential-feedback article has a concrete answer for the class
\eqref{eq:Q}. The coefficient is governed by one explicit scalar saddle,
a Gaussian large action, and a Poisson cloud of twos. The cloud supplies a
factor \(e^{r_n^2}\) that would be missed by retaining only a large action
and unit actions. The unit-action feedback supplies the additional
\(e^{\lambda_1r_n^2}\), explaining both finite-prefix universality and
its sharp exception.

The resulting equivalent identifies a hierarchy of perturbation scales,
strengthens Borel growth from double-logarithmic precision to a full
multiplicative formula, and determines the least-term scale. These results
are proved for a precisely specified model and tail class. They do not
settle unrestricted transseries summability, signed reversion, or the
optimal remainder of a sectorial realization. Those remaining questions
now have more precise parameters and sharper candidate scales.

\appendix
\section{Notation and proof-audit checkpoints}
\label{app:audit}
\begin{center}
\small
\begin{tabular}{@{}p{0.22\textwidth}p{0.66\textwidth}@{}}
\toprule
Symbol & Meaning\\
\midrule
\(u_n\) & Coefficient of the formal fixed point, without factorial scaling.\\
\(b\) & The fixed number \(1+\lambda_1=u_2\).\\
\(b_n\) & The separate computational quantity \(n!u_n\).\\
\(K,J,D_2\) & Number of factors, largest action, number of twos.\\
\(R_n\) & The excess-action coordinate \(n-K+1\).\\
\(r_n\) & Positive solution of \(r_n(r_n+1)e^{2r_n}=n\).\\
\(x_n,m_n\) & Saddle large action and remaining unit mass,
\(n/(r_n+1)\) and \(nr_n/(r_n+1)\).\\
\(D(r),\sigma_n^2\) & \(2r^2+4r+1\) and \(m_n/D(r_n)\).\\
\(A_J\) & Row Poisson mean \((n-J)^2/\lambda_J\).\\
\(s\) & Borel saddle parameter, \(s(s+1)=t\).\\
\(R\) & Least-term parameter \(\tfrac12\log(1/t)\), not \(R_n\).\\
\bottomrule
\end{tabular}
\end{center}

An independent proof audit can focus on five points. First, verify the
marking factor in \eqref{eq:marked}: multiple largest parts must be
\emph{overcounted}, not treated as exact. Second, keep the factor
\(e^{\mu_2}\); discarding the entropy of the twos loses a diverging
multiplicative factor. Third, check that the unspecified tail in
\eqref{eq:Q} yields a \emph{uniform} \(o(1)\) in
\eqref{eq:tail-expansion}, without assuming smoothness. Fourth, distinguish
total variation from rescaled pointwise local limits; the extra bound in
\cref{cor:local} is necessary. Fifth, keep the positive-axis Borel estimate
and the least-term minimization separate from an analytic remainder theorem.

\begin{thebibliography}{9}
\raggedright

\bibitem{repo-classification}
ProveIt research archive.
\emph{A Sharp Regularity Classification for Countable Exponential-Feedback
Transseries}. Research draft, 29 September 2026.
Repository maintained by Vladimir Reshetnikov; inspected at commit
\texttt{04473354a0f3edff2d6c365caf26170ca6b88735}.
\href{https://github.com/VladimirReshetnikov/ProveIt/blob/04473354a0f3edff2d6c365caf26170ca6b88735/Analysis/Transseries/docs/series-and-transseries/Exponential_Feedback_Regularity_Classification/article.tex}{Pinned article source}.
The main theorems occur in source lines 140--330; the positive envelope in
lines 430--690; the finer-asymptotics question in the research-question
subsection within lines 1430--1480.

\bibitem{repo-volume}
ProveIt research archive.
\emph{Transseries: the polynomial--logarithmic calculus and its inversions}.
Canonical volume and status README, inspected at the same pinned commit.
\href{https://github.com/VladimirReshetnikov/ProveIt/blob/04473354a0f3edff2d6c365caf26170ca6b88735/Analysis/Transseries/docs/series-and-transseries/Transseries_And_Inversion/README.md}{Pinned volume README and formalization-status map}.

\bibitem{gessel}
Ira M.~Gessel.
\emph{Lagrange Inversion}.
\emph{Journal of Combinatorial Theory, Series A} \textbf{144} (2016), 212--249.
\href{https://arxiv.org/abs/1609.05988}{arXiv:1609.05988}.

\bibitem{janson}
Svante Janson.
\emph{Simply generated trees, conditioned Galton--Watson trees, random
allocations and condensation}.
\emph{Probability Surveys} \textbf{9} (2012), 103--252.
\href{https://arxiv.org/abs/1112.0510}{arXiv:1112.0510}.

\bibitem{sauzin}
David Sauzin.
\emph{Introduction to 1-summability and resurgence}. 2014.
\href{https://arxiv.org/abs/1405.0356}{arXiv:1405.0356}.

% ed. (2026-09-29): bibliography entry added for the editorial note after the abstract.
\bibitem{ed:finitecore}
[Editorial addition, ProveIt, 2026-09-29.]
\emph{Finite-Core Universality and Sharp Large Order for Countable
Exponential-Feedback Transseries}. Research draft, 29 September 2026.
Directory \path{Finite_Core_Universality_Exponential_Feedback/} of
\path{Analysis/Transseries/docs/series-and-transseries/}; filed in batch 47,
after the snapshot inspected by this article.
\end{thebibliography}
\end{document}
